%=====================================================================
\chapter{Responsible and Ethical Machine Learning}\label{ch:ethics}
%=====================================================================
\locator{6}{Part VI --- Responsible Practice}

\begin{outcomes}
By the end of this chapter you will be able to:
\begin{tight}
  \item \textbf{Identify} where bias can enter a machine learning system, from data
        collection to deployment.
  \item \textbf{Compute} and \textbf{compare} group fairness measures --- selection
        rate, equal opportunity and predictive parity --- from confusion matrices.
  \item \textbf{Explain} why the common fairness criteria cannot all be satisfied at
        once when base rates differ.
  \item \textbf{Discuss} privacy, transparency, accountability and security as
        obligations of anyone who deploys a model.
  \item \textbf{Apply} a responsible-ML checklist to a project, and \textbf{write} a
        model card.
\end{tight}
\end{outcomes}

\begin{prereq}
\chref{evaluation} (confusion matrices, precision, recall, thresholds) and
\chref{workflow} (data and features). This chapter delivers the HEC 2025 outcome
``discuss ethical considerations while applying ML to societal problems''.
\end{prereq}

\begin{keyterms}
algorithmic bias $\bullet$ historical, representation and measurement bias
$\bullet$ proxy variable $\bullet$ feedback loop $\bullet$ protected attribute
$\bullet$ demographic parity $\bullet$ disparate impact (four-fifths rule)
$\bullet$ equal opportunity $\bullet$ equalised odds $\bullet$ predictive parity
$\bullet$ impossibility of fairness $\bullet$ privacy $\bullet$ re-identification
$\bullet$ explainability $\bullet$ model card $\bullet$ accountability $\bullet$
adversarial example $\bullet$ data poisoning
\end{keyterms}

\section{Why a Machine Learning Course Has an Ethics Chapter}

A model is a decision procedure applied at scale. A biased hiring manager affects the
candidates they interview; a biased screening model affects every candidate, the same
way, every year, and looks objective while doing it. The earlier chapters taught how
to make a model accurate. This one asks the questions an accurate model does not
answer: accurate for whom, measured how, used to decide what, and with what recourse
for the people it gets wrong? None of these questions is a matter of opinion to be
left to someone else. Each has technical content, and the engineer who built the model
is often the only person who can answer it.

\section{Where Bias Enters}

\dsfigh{%
\begin{tikzpicture}[font=\scriptsize,
  st/.style={rounded corners=3pt, draw=primary, line width=0.9pt, fill=primary!8,
             text=primary!80!black, font=\scriptsize\bfseries, minimum width=2.3cm,
             minimum height=0.8cm, align=center},
  bi/.style={rounded corners=2pt, draw=accent!70, fill=accent!6, text=accent!80!black,
             font=\tiny, text width=2.3cm, align=center, inner sep=3pt},
  a/.style={-{Stealth[length=2mm]}, gray!60, line width=0.9pt}]
\node[st] (s1) at (0,0) {World};
\node[st] (s2) at (2.8,0) {Data collected};
\node[st] (s3) at (5.6,0) {Labels \&\\features};
\node[st] (s4) at (8.4,0) {Model \&\\objective};
\node[st] (s5) at (11.2,0) {Deployment};
\foreach \a/\b in {s1/s2,s2/s3,s3/s4,s4/s5}{\draw[a] (\a) -- (\b);}
\draw[a, accent!70, dashed] (s5.north) .. controls +(0,1.2) and +(0,1.2) ..
     node[above, font=\tiny, text=accent!80!black] {feedback loop: decisions change the world, which becomes next year's data} (s1.north);
\node[bi, anchor=north] at (0,-0.6) {\textbf{historical}: the past was unfair, and the data records it};
\node[bi, anchor=north] at (2.8,-0.6) {\textbf{representation}: some groups barely appear};
\node[bi, anchor=north] at (5.6,-0.6) {\textbf{measurement}: the label is a proxy (tickets closed for productivity, past promotions for talent)};
\node[bi, anchor=north] at (8.4,-0.6) {\textbf{aggregation}: one model, one threshold for groups that differ};
\node[bi, anchor=north] at (11.2,-0.6) {\textbf{use}: a score built for one purpose applied to another};
\end{tikzpicture}}%
{Five places bias enters a machine learning system. Removing a protected attribute
addresses none of them.}{bias-sources}

Two points from \dsref{bias-sources} catch almost every team. First, a model does not
need to see a \term{protected attribute} (gender, religion, ethnicity, region) to
discriminate on it. Other features act as \term{proxy variables}: university, home
district, name, even a gap in a CV. In this chapter's lab, three ordinary features
recover group membership with 97\% accuracy. Second, \term{feedback loops}: a model
that gives the most visible projects to the developers it rates highest hands them the
achievements that make next year's model rate them highest again. The model's
prediction makes itself come true.

\section{Measuring Fairness}

Fairness cannot be measured without choosing what to compare. Four group criteria
are in common use, each a comparison of confusion-matrix quantities between groups:

\rowcolors{2}{white}{lightbg}
\begin{longtable}{@{}L{3.3cm}L{4.6cm}L{7.1cm}@{}}
\toprule
\rowcolor{primary!15}
\textbf{Criterion} & \textbf{Equal across groups} & \textbf{Meaning} \\
\midrule
\endhead
Demographic parity & Selection rate $P(\hat{y} = 1)$ & Each group is selected in the same
proportion. The \term{four-fifths rule} of US employment law flags a ratio below 0.8 as
\term{disparate impact} \\
Equal opportunity & True positive rate (recall) & Qualified people are selected at the
same rate in every group \\
Equalised odds & TPR and FPR & As above, and unqualified people are wrongly selected at
the same rate \\
Predictive parity & Precision $P(y = 1 \mid \hat{y} = 1)$ & A positive prediction means
the same in every group \\
\bottomrule
\end{longtable}

\begin{worked}[Auditing a Hiring Screen]
A software company's model screens 1{,}000 applicants for developer roles from each of
two groups. Group A has 400 truly qualified applicants; group B, whose members had fewer
chances to build a strong record (fewer internships, less access to paid courses), has
200.

\smallskip
\begin{center}
\begin{tabular}{@{}lrrrrr@{}}
\toprule
\textbf{Group} & \textbf{TP} & \textbf{FN} & \textbf{FP} & \textbf{TN} & \textbf{Selected} \\
\midrule
A & 320 & 80 & 60 & 540 & 380 \\
B & 120 & 80 & 40 & 760 & 160 \\
\bottomrule
\end{tabular}
\end{center}
\smallskip
\begin{tight}
  \item \textbf{Selection rate:} A $380/1000 = 38\%$; B $160/1000 = 16\%$. Ratio
        $0.16/0.38 = \mathbf{0.42}$ --- far below 0.8: disparate impact.
  \item \textbf{True positive rate:} A $320/400 = 0.80$; B $120/200 = \mathbf{0.60}$. A
        qualified B applicant is rejected 40\% of the time, twice as often as a
        qualified A applicant (20\%). Equal opportunity fails.
  \item \textbf{False positive rate:} A $60/600 = 0.10$; B $40/800 = 0.05$.
  \item \textbf{Precision:} A $320/380 = 0.84$; B $120/160 = 0.75$.
\end{tight}
\textbf{Reading it.} Part of the gap in selection reflects a real difference in base
rates (40\% against 20\% qualified); but the gap in \emph{true positive rate} does not ---
it is the model treating equally qualified people differently. That is the finding to
act on first.
\end{worked}

\dsfig{%
\begin{tikzpicture}
\begin{axis}[width=10.4cm, height=5.0cm, ybar, bar width=0.42cm, axis lines=left,
  ymin=0, ymax=1.0, ylabel={rate}, symbolic x coords={selection,TPR,FPR,precision},
  xtick=data, tick label style={font=\tiny}, label style={font=\scriptsize},
  enlarge x limits=0.18, nodes near coords, nodes near coords style={font=\tiny},
  legend style={font=\tiny, draw=gray!40, at={(0.02,0.97)}, anchor=north west},
  point meta=explicit symbolic]
\addplot[fill=secondary!55, draw=secondary] coordinates
  {(selection,0.38)[0.38] (TPR,0.80)[0.80] (FPR,0.10)[0.10] (precision,0.84)[0.84]};
\addlegendentry{group A}
\addplot[fill=accent!50, draw=accent] coordinates
  {(selection,0.16)[0.16] (TPR,0.60)[0.60] (FPR,0.05)[0.05] (precision,0.75)[0.75]};
\addlegendentry{group B}
\end{axis}
\end{tikzpicture}}%
{The worked audit. Every criterion shows a gap, but they mean different things: the
selection gap partly reflects different base rates; the TPR gap is unequal treatment of
equally qualified applicants.}{fairness-bars}

\begin{alertbox}[The Impossibility of Fairness]
When two groups have different base rates, no classifier that makes any errors can
satisfy equal TPR, equal FPR and equal precision at the same time (Chouldechova, 2017;
Kleinberg, Mullainathan and Raghavan, 2016). Choosing a fairness criterion is
therefore a \emph{decision}, not a calculation: equal opportunity protects qualified
applicants; predictive parity protects the meaning of a positive score. In the lab,
lowering group B's threshold to 0.26 equalises the TPR at 0.93 --- and B's precision
falls from 0.83 to 0.63. Say which criterion you chose, why, and who pays for it.
\end{alertbox}

\section{Privacy}

Training data often describes people, and a model can leak it. Removing names is not
anonymisation: a small number of attributes --- postcode, date of birth and gender ---
identifies most individuals in a population, and linking a ``de-identified'' dataset to a
public one \term{re-identifies} them. Models themselves can memorise and reveal rare
training examples. Good practice: collect only what the task needs; keep identifiers
out of features; aggregate or add noise where exact values are not needed
(\emph{differential privacy} makes that noise mathematically rigorous); and control
who can query a model trained on sensitive data. Legal duties vary by jurisdiction ---
Pakistan's Prevention of Electronic Crimes Act 2016, the Personal Data Protection Bill
drafted by the Ministry of IT and Telecommunication, the EU's GDPR for data about
European residents --- and a project handling personal data should identify which apply
before collecting anything.

\section{Transparency and Explainability}

People affected by a decision are entitled to an explanation they can act on, and the
people deploying a model need to know when it is relying on something it should not.
Three levels help:
\begin{tight}
  \item \textbf{Interpretable models} where the stakes justify the accuracy cost:
        logistic regression's odds ratios (\chref{logreg}), a shallow tree's rules
        (\chref{trees}).
  \item \textbf{Global explanations} of complex models: permutation importance ---
        how much does performance drop when one feature is shuffled? --- and partial
        dependence.
  \item \textbf{Local explanations} of a single prediction: which features pushed
        \emph{this} applicant's score down (SHAP values and LIME are the common tools).
\end{tight}
An explanation of what the model \emph{uses} is not an explanation of what
\emph{causes} the outcome (\chref{linreg}); say which one you are giving.

\begin{definitionbox}[Model Card]
A \term{model card} is a short document published with a model: what it is for, and
what it must not be used for; the data it was trained on; its performance overall
\emph{and for each relevant group}; the fairness criterion chosen and the result; known
limitations; and who is responsible for it. The project of \chref{project} must include
one.
\end{definitionbox}

\section{Accountability and Security}

\textbf{Accountability} means that a person, not ``the algorithm'', is responsible for
a decision, and that the people affected can contest it. Keep a human in the loop for
consequential decisions; log what the model predicted and why; and provide a route for
appeal that can actually overturn the prediction.

\textbf{Security} is an ethical issue too, because a model that can be manipulated will
be. \term{Adversarial examples} --- inputs changed slightly and deliberately --- flip a
model's output (\chref{ann}). \term{Data poisoning} inserts crafted examples into the
training data, a particular risk for models retrained on user-supplied data and for
self-training (\chref{semisup}). A model deployed where someone benefits from fooling
it must be tested against someone trying.

\rowcolors{2}{white}{lightbg}
\begin{longtable}{@{}L{3.6cm}L{11.4cm}@{}}
\toprule
\rowcolor{primary!15}
\textbf{Stage} & \textbf{Questions to answer before deployment} \\
\midrule
\endhead
Purpose & What decision will this support? Who is affected? Would a simpler rule, or no
model, be better? \\
Data & Who is missing or under-represented? What is the label really measuring? Was it
collected with consent, and may it be used this way? \\
Features & Could any feature be a proxy for a protected attribute? Does any exist only
because of the outcome (\chref{workflow})? \\
Evaluation & Performance for each group, not only overall. Which fairness criterion,
and why? \\
Deployment & Who reviews the model's decisions? How can they be contested? How will
drift and feedback loops be detected? \\
Security & Who benefits from fooling it? Has it been tested against them? \\
\bottomrule
\end{longtable}

\begin{pitfall}
\begin{tight}
  \item \textbf{``We removed the protected attribute, so the model is fair.''}
        Proxies carry it back in.
  \item \textbf{Reporting only overall accuracy.} A model can be accurate on
        average and poor for a minority group; report per group.
  \item \textbf{Picking the fairness metric the model happens to pass.} Choose it
        from the harm you want to prevent, before looking.
  \item \textbf{Treating anonymised data as safe.} Linking attributes re-identifies
        people.
  \item \textbf{Deploying without a way to be wrong.} Every model errs; the question is
        what happens to the person it errs about.
\end{tight}
\end{pitfall}

%---------------------------------------------------------------------
\begin{lab}[Auditing a Model for Fairness]
A synthetic hiring dataset in which group B had fewer chances to build a strong
record. The group is never given to the model as a feature.
\begin{lstlisting}[language=Python]
import numpy as np
from sklearn.linear_model import LogisticRegression
from sklearn.model_selection import train_test_split, cross_val_score

# --- a synthetic hiring dataset with a protected group -------------------
rng = np.random.default_rng(0)
n = 6000
group = rng.integers(0, 2, n)                     # 0 = group A, 1 = group B
# group B had fewer internships and projects to show (historical bias)
record = rng.normal(0.6 - 0.25 * group, 0.15, n)  # CV: internships, projects
aptitude = rng.normal(0.5, 0.15, n)               # coding test: equal
district = group + rng.normal(0, 0.3, n)          # a proxy for the group
qualified = (0.5 * record + 0.5 * aptitude
             + rng.normal(0, 0.05, n) > 0.5).astype(int)
X = np.column_stack([record, aptitude, district])
X_tr, X_te, y_tr, y_te, g_tr, g_te = train_test_split(
    X, qualified, group, test_size=0.5, random_state=0)


def report(name, pred):
    print(name)
    for g, label in ((0, "A"), (1, "B")):
        m = g_te == g
        sel = pred[m].mean()                               # selection rate
        tpr = pred[m][y_te[m] == 1].mean()                 # recall
        fpr = pred[m][y_te[m] == 0].mean()
        prec = y_te[m][pred[m] == 1].mean()
        print(f"  group {label}: base rate {y_te[m].mean():.2f}  "
              f"selected {sel:.2f}  TPR {tpr:.2f}  FPR {fpr:.2f}  "
              f"precision {prec:.2f}")


# --- 1. the model, group column excluded (it was never a feature) --------
model = LogisticRegression().fit(X_tr, y_tr)
p = model.predict_proba(X_te)[:, 1]
report("threshold 0.5 for everyone:", (p > 0.5).astype(int))

# --- 2. removing the group column does not remove the group --------------
acc = cross_val_score(LogisticRegression(), X, group, cv=5).mean()
print(f"the features recover group membership with accuracy {acc:.2f}")

# --- 3. equalise the true positive rate with a group-specific threshold --
t_b = 0.5
tpr_a = (p[(g_te == 0) & (y_te == 1)] > 0.5).mean()
while (p[(g_te == 1) & (y_te == 1)] > t_b).mean() < tpr_a:
    t_b -= 0.01
pred = np.where(g_te == 1, p > t_b, p > 0.5).astype(int)
report(f"threshold 0.5 for A, {t_b:.2f} for B (equal opportunity):", pred)
\end{lstlisting}
Equalising the TPR costs group B precision and raises its false positive rate --- the
impossibility result in action. Write two sentences, as if for the hiring committee,
explaining the trade-off you would recommend.
\end{lab}

%---------------------------------------------------------------------
\begin{checkpoint}
\begin{tightnum}
  \item A hiring model selects 45\% of one group and 30\% of another. Does it pass the
        four-fifths rule?
  \item Why does removing a protected attribute from the features not guarantee a fair
        model?
  \item Which fairness criterion protects qualified applicants from being rejected,
        and which confusion-matrix quantity does it compare?
\end{tightnum}
\end{checkpoint}

%---------------------------------------------------------------------
\begin{chaptersummary}
\begin{tight}
  \item Bias enters through history, representation, measurement, aggregation and
        use, and feedback loops can amplify it; proxies reintroduce removed
        attributes.
  \item Group fairness compares confusion-matrix rates across groups: selection
        (demographic parity), TPR (equal opportunity), TPR and FPR (equalised odds),
        precision (predictive parity).
  \item With different base rates these criteria conflict; choosing one is a
        decision to be justified, not computed.
  \item Protect privacy by minimising and aggregating data; explain models at the
        right level; keep a human accountable and a route to contest.
  \item Document every deployed model with a model card, including per-group
        performance, and test it against people with a reason to fool it.
\end{tight}
\end{chaptersummary}

\begin{reviewq}
  \item \textbf{(MCQ)} Equal opportunity requires equal: (a)~selection rates
        (b)~true positive rates (c)~precision (d)~accuracy.
  \item \textbf{(MCQ)} Using ``tickets closed per month'' as the label for ``good
        developer'' is an example of: (a)~representation bias (b)~measurement bias
        (c)~a feedback loop only (d)~data poisoning.
  \item \textbf{(MCQ)} A selection-rate ratio of 0.42 between two groups: (a)~passes the
        four-fifths rule (b)~fails it (c)~proves intent to discriminate (d)~says nothing.
  \item \textbf{(Short)} Explain a feedback loop with an example of your own from IT
        project management.
  \item \textbf{(Short)} Why can equal opportunity and predictive parity not both hold
        when base rates differ?
  \item \textbf{(Short)} List the sections of a model card and justify one of them.
  \item \textbf{(Applied)} For the worked audit, find how many more group-B qualified
        applicants must be selected to equalise the TPR at 0.80, and compute B's new
        selection rate if no additional false positives are admitted.
  \item \textbf{(Applied)} Choose one model from Chapters 5--20 and write its model
        card for a realistic deployment, using the checklist of this chapter.
\end{reviewq}
